\section{研究任务与数据来源}

本项目关注的问题是：当大语言模型面对一个关于“电动轮椅能否在楼梯转角处完成 $90^\circ$ 转弯”的空间推理问题时，它的最终判断是否稳定、是否受表达层级和无关扰动影响，以及不同模型之间是否呈现可挖掘的失效模式。这里的核心对象不是单个几何题答案，而是模型在受控 Prompt 条件下产生的行为记录。

上游 Prompt 生成部分已经把一个空间场景表示为五元组
\begin{equation}
\mathbf{x}=(W,D,w,L,R),
\end{equation}
其中 $W$ 是楼梯通道净宽，$D$ 是平台有效进深，$w$ 和 $L$ 分别表示轮椅宽度与长度，$R$ 是厂家或题干给出的保守最小转弯半径，单位均为厘米。为了给后续评价提供一致参照，数据集中定义几何余量
\begin{equation}
\kappa(\mathbf{x})=\min(W-w,\ D-R).
\label{eq:clearance-margin}
\end{equation}
若 $\kappa\geq 3$，参考标签记为 \code{can_pass}；否则记为 \code{cannot_pass}。这些标签不会暴露给模型，只在后续特征构建和准确性分析时使用。

\begin{figure}[t]
\centering
\begin{tikzpicture}[scale=0.78, every node/.style={font=\scriptsize}]
  \fill[dmgray] (-0.25,-0.25) rectangle (5.15,3.65);
  \fill[white] (0,0) rectangle (4.8,1.15);
  \fill[white] (3.65,0) rectangle (4.8,3.4);
  \draw[black!55, line width=0.45pt] (0,1.15) -- (3.65,1.15) -- (3.65,3.4);
  \draw[black!55, line width=0.45pt] (0,0) -- (4.8,0) -- (4.8,3.4);
  \draw[dmorange, dashed, line width=0.7pt] (2.72,1.15) arc[start angle=180,end angle=90,radius=0.93];
  \node[dmnote, dmorange] at (3.15,1.78) {$R$};
  \draw[fill=dmgreen!18, draw=dmgreen!80, rounded corners=1.5pt, rotate around={8:(2.65,0.55)}]
    (2.22,0.26) rectangle (3.08,0.84);
  \draw[<->, black!65] (2.22,0.08) -- (3.08,0.08)
    node[midway,below] {$L$};
  \draw[<->, black!65] (3.25,0.26) -- (3.25,0.84)
    node[midway,right] {$w$};
  \draw[<->, black!65] (-0.1,0) -- (-0.1,1.15)
    node[midway,left] {$W$};
  \draw[<->, black!65] (3.65,3.55) -- (4.8,3.55)
    node[midway,above] {$D$};
  \node[dmnote] at (2.0,1.42) {楼梯转角平台};
  \node[dmnote] at (4.25,2.55) {转向通道};
\end{tikzpicture}
\caption{轮椅转弯空间推理任务的抽象几何示意。实验并不追求真实工程仿真，而是固定 $W,D,w,L,R$ 等关键变量，用受控 Prompt 观察模型如何利用或误用这些空间信息。}
\label{fig:task-geometry}
\end{figure}


每条 Prompt 可以抽象为
\begin{equation}
p=f(\mathbf{x},\ell,n,g),
\end{equation}
其中 $\ell\in\{\mathrm{L1},\mathrm{L2},\mathrm{L3}\}$ 表示描述层级，$n$ 表示扰动条件，$g$ 表示场景生成策略。模型 $m$ 在第 $r$ 次重复采样中的回复记为 $a_{m,r}(p)$。这一记号使后续分析能够直接按策略、层级、扰动、模型和重复编号进行分组。
